% Generated by study2b/scripts/make_paper_example.py from the committed data.
% Do not edit by hand: change the script and re-run it.
\begin{figure}
  \footnotesize
  \renewcommand{\arraystretch}{1.15}
  \begin{tabular}{@{}>{\raggedright\arraybackslash}p{0.25\linewidth}p{0.71\linewidth}@{}}
    \toprule
    PURE source & The scheduler shall support up to 3 rate groups per vehicle mode. \\
    Input (what the model sees) & \emph{``So for the scheduling piece, I think it needs to handle a handful of rate groups per vehicle mode --- three feels like plenty, that's the ceiling we'd want it to work with, don't need more than that realistically.''} \\
    Admissible & Ubiquitous only (both annotators) \\
    \midrule
    Base prompt & WHERE the scheduling feature is included, the system shall support scheduling for up to three rate groups per vehicle mode.\newline \textbf{Checker: compliant. Pattern fit: no.} \\
    Base prompt with examples & The scheduling component shall support up to three rate groups per vehicle mode.\newline \textbf{Checker: compliant. Pattern fit: yes.} \\
    \bottomrule
  \end{tabular}
  \caption{Comparison of the EARS patterns chosen for one input under two prompts.}
  \Description{A table for one input. It gives the PURE source requirement, the informal input the model saw, and the admissible pattern (Ubiquitous only). It then shows two outputs: with the base prompt, a requirement opening with a WHERE clause, which passes the checker but does not fit; with worked examples, an unconditional requirement, which passes the checker and fits.}
  \label{fig:example}
\end{figure}
